\chapter{Getting grammaticality wrong}\label{ch:getting-grammaticality-wrong}

A book about grammaticality has to make room for a hard fact: experts can be wrong about it. Not just careless, not merely prescriptive, but wrong in the ordinary evidential sense. They can mistake the silence of their own experience for the silence of the language.

\section{The diachronic context}

% TODO. Keep the existing material below — Old English þe, the Middle English
% emergence of relative who, the þe his construction as the precursor to
% whose. Sets up the modern reveal.
% Existing content: ~25 lines; expand if useful.

% TODO: weaken claim — OE had *se/seo/þæt* and *se þe* relative constructions; current claim ignores them
Old English didn't have relative pronouns the way we have them. It just used \mention{þe} to mark subordinate clauses of various kinds, including relative clauses.

% TODO: cite source edition (kyninga)
% TODO: critical — verify the kyninga example against an actual OE source. The genitive plural marking suggests *þonne* is the comparative ("than"), not temporal ("Then"); current translation as a relative-clause-modifying-subject may be wrong. Bert-Remijsen-risk-level OE-data error per the proofread agent.
\ea \label{ex:kyninga}
    \gll þonne ealra oþra kyninga þe in middangearde æfre wæron \dots\\
        Then all other kings that in Middle-Earth ever were \dots\\
    \glt `Then all the kings who were ever on earth \dots'
\z

There was \mention{hwa} `who', but it was mainly interrogative. It had a special construction meaning roughly `whoever', as in \mention{whoever enters remains}, but the \mention{þe} in (\ref{ex:kyninga}) couldn't have been replaced with \mention{hwa}. It was only around 1325, in the midst of Middle English, that we start to see a real relative \mention{who} in the way we know it today.

% TODO: cite source edition (doȝter wo, plausibly Robert of Gloucester c. 1325)
\ea \label{ex:doȝter-wo}
    \gll He nadde bote an doȝter wo miȝte is eir be.\\
        He hadn't but one daughter who might his heir be.\\
    \glt `He had no one but a daughter who could be his heir.'
\z

Similarly, there wasn't a relative \mention{whose} in Old English. Instead you get constructions like \mention{þe his} `that his' in (\ref{ex:þe-his}).

% TODO: cite source edition (Eaðig bið, plausibly the OE Psalter Ps 39:5/40:4)
\ea \label{ex:þe-his}
    \gll Eaðig bið se wer, þe his tohopa bið to Drihtne\\
        Blessed be the man that his hope is in the Lord.\\
    \glt `Blessed be the man whose hope is in the Lord'
\z

\section{The Hankamer-Postal moment}

The article \enquote{Whose gorilla?}, which appeared in \textit{Linguistic Inquiry} in \citeyear{Hankamer1973}, was a \enquote{squib}~-- a short publication that presented an interesting piece of data without proposing an explanation. Not even running a full page, its punch line was (\ref{ex:gorilla}):

\ea \label{ex:gorilla}
 \textit{~My gorilla is over there drinking punch.} \\
 *\textit{The guy whose you saw banging at the window is over there watering}\\~\textit{the rubber tree.}
\z
They had fun in those days.

The authors of \enquote{Whose gorilla?} were Jorge Hankamer and Paul Postal. \enquote{Postal was a venerable force, \dots~ second in reputation only to Chomsky in the prized domain of syntax} \citep[127--128]{Harris2021}. He was also an outsider, spending most of his academic life at IBM instead of at a university, and something of a polemicist. His most famous paper was entitled \enquote{The best theory}, and it supported a position called generative semantics that Chomsky would later deride as \enquote{the worst theory, in fact as the worst possible theory \dots~the worst imaginable theory} \citep[129]{Harris2021}. This despite the irony that \enquote{the Chomsky of the 21st century reads like a reincarnation of the Postal of the 20th} \citep[418]{Newmeyer2003}.

As a graduate student, Postal was so sure of himself that he directed his own thesis. When he submitted it, his advisor wrote him a note saying \enquote{he was too busy with his work to read it for two years} \citep{PostalInterview2020}. This mixture of brilliance and confidence made him one of the key combatants in the so-called Linguistics Wars of the late `60s and early `70s, a truly rancorous battle between the generative semanticists on one side and Chomsky and his supporters on the other.

In 1973, Hankamer had just completed his PhD at Yale and taken up a position at Harvard. His dissertation, on deletion and ellipsis in Turkish, had already established him as a rising star in the field. Like Postal, he was on the outs with Chomsky despite being far less bombastic and pugilistic.

In their article \enquote{Whose gorilla?}, Hankamer and Postal drew attention to an apparent gap in the pronoun system, specifically claiming that the word \mention{whose} is defective.

In linguistic terminology, a \textsc{defective} word is one that lacks one or more of the expected forms. This is why Hankamer and Postal have starred that sentence, to draw attention to the \mention{whose} that should~-- but seemingly couldn't~-- be there.

What does it mean to lack a form? \mention{Beware}, for example is missing every form but one; it's highly defective. It has a bare form \mention{beware of the dog}, but no present-tense *\mention{she bewares the dog}, no past tense *\mention{She bewore the dog}, and no participles *\mention{I've never beworn the dog}.

This missing form~-- the purported hole in the \mention{whose} paradigm~-- should appear in the bottom row of Table \ref{tab:pronouns}, right where the question mark is.
%, though many linguists prefer the term \textsc{genitive}, which suggests a wider set of relationships than simple possession. When you speak of \uline{your} ancestors, \uline{your} birth, or \uline{your} situation~-- \textit{your} being a possessive/genitive pronoun~-- clearly you do not mean that you possess these things. A simple way to remember this is that the \uline{gen}itive suggests a \uline{gen}eral connection between things. Terminological issues aside, \textit{whose gorilla} evokes some kind of relationship with the gorilla.

\begin{table}[h]
\centering
\begin{tabular}{lccccc}
\hline
\textbf{Type} & \textbf{Nominative} & \textbf{Accusative} & \textbf{Dep Gen} & \textbf{Ind Gen} & \textbf{Reflexive} \\
\hline
\textbf{Personal} & \textit{I} & \textit{me} & \textit{my} & \textit{mine} & \textit{myself} \\
 & \textit{you} & \textit{you} & \textit{your} & \textit{yours} & \textit{yourself} \\
 & \textit{he} & \textit{him} & \textit{his} & \textit{his} & \textit{himself} \\
 & \textit{she} & \textit{her} & \textit{her} & \textit{hers} & \textit{herself} \\
 & \textit{it} & \textit{it} & \textit{its} & - & \textit{itself} \\
 & \textit{we} & \textit{us} & \textit{our} & \textit{ours} & \textit{ourselves} \\
 & \textit{they} & \textit{them} & \textit{their} & \textit{theirs} & \textit{themselves} \\
 & \textit{one} & \textit{one} & \textit{one's} & \textit{one's} & \textit{oneself} \\
\textbf{Interrog} & \textit{who} & \textit{whom} & \textit{whose} & \textit{whose} & - \\
\textbf{Relative} & \textit{who} & \textit{whom} & \textit{whose} & ? & - \\
\hline
\end{tabular}
\caption{Personal, Interrogative, and Relative Pronouns}
\label{tab:pronouns}
\end{table}

\mention{Whose} is often known as a \textsc{genitive pronoun}, AKA \enquote{possessive pronoun}. The genitive pronouns come in two kinds: those that are \textsc{dependent} on a noun: \mention{\uline{my} life}, \mention{\uline{your} name}, \mention{\uline{her} eyes}, and those that are \textsc{independent}: \mention{mine}, \mention{yours}, \mention{hers}, etc. Perhaps confusingly, dependent and independent forms can share the same spelling and pronunciation, as in \mention{it's \uline{his} job} (dependent) or \mention{the job is \uline{his}} (independent). Such is the case with \mention{whose}: You can ask \mention{whose gorilla is this} or \mention{whose is it}, at least when you use \mention{whose} to ask a question.

The \mention{whose} of questions is an \textsc{interrogative pronoun}, but there's another \mention{whose}, as in \mention{that's the person whose name is on the door}. This is the {relative pronoun}. Relative \mention{whose} isn't interrogating anything; it's \enquote{standing in for} \mention{the person's}. And the \mention{whose} in (\ref{ex:gorilla})~-- the one that Hankamer and Postal marked as ungrammatical~-- is a relative pronoun. Putting it all together, then, the \mention{whose} at issue is this independent genitive relative pronoun.

The point Hankamer and Postal were making was that, where most pronouns have paired dependent and independent genitives, \mention{who} feels like it's missing the independent one. All the other combinations in (\ref{ex:whose-good}) are possible. Why not (d)?

\ea \label{ex:whose-good}
 \ea[]{\textit{Whose gorilla is this?}\hfill[dependent interrogative]}
 \ex[]{\textit{Whose is this?}\hfill[independent interrogative]}
 \ex[]{\textit{someone whose gorilla had left}\hfill[dependent relative]}
 \ex[*]{\textit{someone whose had left}\hfill[independent relative]}
 \z
\z

They didn't know, and that's why they published \enquote{Whose gorilla} as a squib. Generally, squibs don't get cited much, and this one, true to form, has been mostly ignored. In 2002 the \textit{Cambridge grammar of the English language} agreed with its judgment, marking (\ref{ex:CGEL-whose}) with an asterisk. And, nearly 20 years after that, almost 30 since the article was published, \citet{Cinque2020a} was still referring to the issue as a \enquote{residual puzzle}.

\ea[*]{\textit{The police are trying to contact the person \ob whose it was\cb.} \hfill(p. 472)}\label{ex:CGEL-whose}
\z

\section{Whose, found in the wild}

How had Hankamer and Postal determined that (\ref{ex:gorilla}) was ungrammatical? Almost certainly by noticing their own metacognitive feelings. Did they search for instances or present relevant examples to a representative sample of English speakers to see how the independent \textit{whose} was perceived? They don't say, but it's very unlikely. Such experimental methods were almost never seen as necessary at the time. A linguist's personal  grammaticality judgments were generally deemed valid and reliable. Rarely were they questioned.

It occurred to me, though, to do something they couldn't have done. In March of 2024, I asked OpenAI's GPT-4 and Google's Gemini Advanced to rate the grammaticality of this \textit{whose} on a scale of 1--7, not because I questioned the feelings of Hankamer and Postal and the Cambridge grammar authors, but because I was curious to see how these large language models would respond. The sentence I gave them was (\ref{ex:whose-wasn't}).

\ea[]{\textit{Mine was working, but I know someone whose wasn't.}} \label{ex:whose-wasn't}
\z

Surprisingly to me, both models rated it a 6, denoting high grammaticality, and cited the slightly unusual use of \textit{whose} as the reason for not awarding a 7. I'd experimented with eliciting grammaticality judgments from large language models before with various constructions, and I'd mostly found their assessments to align surprisingly well with my own. I'd even asked them to produce examples at various points on the scale, and again, they'd fared reasonably well. This is why I was rather surprised at a 6 rating for (\ref{ex:whose-wasn't}).

Intrigued by this mismatch, it occurred to me to ask some non-linguist colleagues and family members what they thought of it. I sent out a few quick texts and within a day, ten of the twelve I'd asked had also given it a 6 or 7. Two gave it a 4. Many didn't mention \textit{whose}, and even after I queried the pronoun specifically, two said it seemed fine to them. To my surprise, the large language models were a better fit for the naive humans than were the linguists. I was intrigued.
% EDITORIAL-HPC: Do not let this become "LLMs + friends outvote experts." The
% HPC-compatible claim is calibration: expert judgments, naive judgments, corpus
% evidence, OED evidence, and LLM outputs are different noisy instruments whose
% errors can be compared.

This issue arose while I was working on \textit{Wikipedia}'s \href{https://en.wikipedia.org/wiki/English_pronouns}{English pronouns} page. There, the other editors and I had reflected the position of Hankamer and Postal and the Cambridge grammar: that independent relative \textit{whose} didn't exist. With these new grammaticality judgments, I started to wonder if I should revise that entry.

\textit{Wikipedia} requires sources, though, and \enquote{Brett's friends and family along with some LLMs} doesn't cut it. It was then that it occurred to me to consult \textit{The Oxford English Dictionary}. Had Hankamer and Postal done so, perhaps they wouldn't have published because they would have found that the lexicographers at the \textit{OED} seem to take the same position as my informants and the large language models: the independent genitive relative \textit{whose} is rare, but grammatical. The \textit{OED} presents examples going back to the 1300s, including those in (\ref{ex:OED-whose}). %even the first edition had these examples.

\ea \label{ex:OED-whose}
 \ea[]{\textit{After þe wille of him \uline{hos} þe werkes bez.} `According to the will of him whose works are.'\hfill[1325]}
 \ex[]{\textit{The man \uline{whose} these are.}\hfill[1611]}
 \ex[]{\textit{The most generous patron \dots~continues to be our \dots~Pittsburg and Allegheny auxiliary, \uline{whose} also is the largest contribution to the rent of our Association's new office.}\hfill[1904]}
 \ex[]{\textit{Everything depends on the person \uline{whose} this administration is.}\hfill[2018]}
 \z
\z

\noindent And (\ref{ex:online-whose}) presents a few more I found online.

\ea \label{ex:online-whose}
 \ea[]{\textit{In California you usually include a copy of the check with the offer. Earnest money despoit (EMD) can be anywhere from 1/2\%-3\% of the sales price. Mine was 3\%, I just helped a person \uline{whose} was 2\%, I've never seen an EMD less than 1/2\% in the past 3 years in California (from Shasta Co. to Orange Co.).}}
 \ex[]{\textit{I have met a person whose SSN was only 10 away from mine (yes only the second to last digit was one higher than mine) and he was born in a different state and was at least 7-10 years older than I was. And I met another person \uline{whose} was kinda close to mine but again we had nothing in common pertaining to our birthday or location of birth/SSN registries.}}
 \ex[]{\textit{Mine is also submitted on June 30 and no updates yet, i know a person \uline{whose} was filed on June 30 as well in Regular mode and he got approval on 8/16. I think we can expect a decision in the next few weeks.}}
 \ex[]{\textit{I'll be scouring eBay for an Alpha keypad. Debating whether to swap it out for the existing one or just use it for programming. The control panel isn't as hard to get to as the person \uline{whose} was installed in the top of a closet.}}
 \ex[]{\textit{Head shaved? Mine wasn't but the guy \uline{whose} was had to give a pubic sample.}}
 \ex[]{\textit{In my opinion, my 50\%/50\% bud/trim mix came in third place. One guy, \uline{whose} was the best in my opinion, he treated his trim exactly like he did his buds.}}
 \ex[]{\textit{When I entered in 1977, I actually had the longest ash, but mine was bent. I came second to a woman \uline{whose} was straight.}}
 \ex[]{\textit{I have vasovagal syncope (a specific thing triggers fainting) and it's not uncommon for a trigger to be urination or defecation. My trigger is medical stuff like needles, which mostly people are sympathetic to. But I did know another woman \uline{whose} was set off by morning urination.}}
 \ex[]{\textit{Meanwhile the original person \uline{whose} this photograph is\dots}}
 \z
\z

Based on this evidence, I revised the \textit{Wikipedia} article to include the two views. But, of the two positions, was it mine whose was right?

This is the problem with our feelings of (un)grammaticality. Hankamer and Postal are top-notch syntacticians, legendary even. And John Payne and Rodney Huddleston, the authors of the chapter of the \textit{Cambridge grammar of the English language} that marks (\ref{ex:CGEL-whose}) as ungrammatical, are among the best English grammarians the world.

On the other hand, there's no theoretical reason to expect relative \textit{whose} to go AWOL, there's the data from the OED and my web searches, there's the intuition of my friends and family, and there's whatever ineffable average over trillions of words of English is represented in the judgments of the large language models.
% EDITORIAL-HPC: "Ineffable average" is evocative but needs methodological
% discipline when this material moves to ch. 14. LLM judgments are useful as
% distribution-sensitive probes, not as direct evidence of grammaticality.

So, we have two seemingly contradictory forces at play. On the one hand, certain established linguistic authorities tell us that they feel the sentence \textit{The person whose wasn't working} is ungrammatical. On the other hand, the feelings of various people, lexicographers and non-linguists alike, backed up by the vast echo chamber of everyday speech reflected in online examples and the judgments of generative AI models, suggest it might be fine.

\section{Why it works when it works}

But language has a way of surprising us. As it turns out, this missing \mention{whose} isn't quite as missing as Hankamer and Postal thought. It's more like a linguistic ivory-billed woodpecker: long thought extinct, occasionally glimpsed, hotly debated, but stubbornly refusing to be definitively ruled out.

Consider this sentence from the \textit{Cambridge Grammar of the English Language}, a book that's about as authoritative as it gets in the world of English grammar:

\ea{\textit{I was going to visit Lucy, }[\textit{a friend of whose had told us of the accident}]\textit{.}}
\z

Here we see our elusive \mention{whose}, standing proudly without a noun in tow. It's right there in the relative clause \textit{a friend of whose had told us of the accident}. No gorillas in sight, but it's the same type of \mention{whose} that Hankamer and Postal claimed didn't exist.

\bigskip

% TODO: rewrite — AI-tic bridge
So what's going on here? How can this \mention{whose} both not exist and exist at the same time? Well, that's where things get really interesting...

The key to understanding this linguistic puzzle lies in the interactions of syntax, semantics, and pragmatics~-- the rules of sentence structure, meaning, and contextual language use. While Hankamer and Postal were correct that this independent relative \mention{whose} is extremely rare, they missed some crucial factors that allow it to exist in very specific circumstances.

First, let's consider the structure. In the sentence from the Cambridge Grammar, \mention{whose} is part of a larger phrase: \textit{a friend of whose}. This structure, called the \enquote{oblique genitive}, seems to provide a comfortable habitat for our rare \mention{whose}. It's like finding an ivory-billed woodpecker not just anywhere, but in a particular type of old-growth forest.

But structure alone doesn't explain everything. After all, we can't say *\textit{The guy of whose is making a scene.} There's more at play here.

The second factor is information structure -- how we package information in a sentence. This is a bit like how we organize items in a suitcase; some things are right on top (easy to access), while others are buried deep inside.

In language, the items that are \enquote{right on top} are said to be topical. They're the main focus of what we're talking about.

In the Cambridge Grammar example:

\ea{\textit{I was going to visit Lucy, }[\textit{a friend of whose had told us of the accident}]\textit{.}}
\z

Both Lucy (the possessor) and her friend (the possessed) are highly topical:

\begin{itemize}
    \item Lucy is important because she's the person being visited.
    \item The friend is important because they're the one who told about the accident.
\end{itemize}

This high topicality of both elements -- the possessor (Lucy) and the possessed (her friend) -- seems to be a necessary condition for our independent \mention{whose} to appear. It's like our rare linguistic bird needs both large tracts of old-growth bottomland hardwood forests and an abundance of large, recently dead trees.

% TODO: rewrite — AI-tic bridge
To understand this better, we need to get up to speed on some fundamental concepts about how language works, particularly pronouns, ellipsis, and their antecedents.

\subsection{Pronouns, ellipsis, and their antecedents}

In language, we often refer back to something we've already mentioned. This backward reference is called anaphora, and the thing we're referring back to is called the antecedent. In the following example, \textit{that guy} refers back to Aden.

\ea \textit{\uline{Aden} just bought a new car. \uline{That guy} sure loves his rides.}\\
    \begin{tikzpicture}[overlay, remember picture]
        \node[align=left,inner sep=0] (text) {\phantom{\textit{\uline{Aden} just bought a new car. \uline{That guy} sure loves cars.}}};
        \draw[->,red,very thick] ([xshift=6.5ex,yshift=0.8ex]text.north east) to [bend left=20] ([xshift=-21ex,yshift=0.8ex]text.north east);
    \end{tikzpicture}\label{ex:Aden1}
\z

Noun phrases like \textit{that guy} can perform this referential function, but pronouns -- small words like \textit{they}, \textit{it}, \textit{he}, and \textit{she} -- are the true specialists in anaphoric reference. For this linguistic magic to work, though, the antecedent must be readily available in the discourse. The surest way to achieve this availability is for the antecedent to be topical or recently mentioned. English provides some clever mechanisms to extend our referential reach: we can distinguish between singular and plural (\textit{she} vs \textit{they}), personal and non-personal (\textit{she} vs \textit{it}), and even masculine and feminine (\textit{he} vs \textit{she}). These distinctions act like linguistic signposts. If we encounter a feminine pronoun, but the most recently mentioned potential antecedent is masculine, it's a cue to search further back in our mental discourse for a more suitable referent. Interestingly, each time we use a pronoun, we're not just referring back -- we're also boosting the prominence of its antecedent in the ongoing discourse. It's as if we're a linguistic juggler, periodically tossing each referent high into the air to keep it in play.

Another way to do this is a bit counter intuitive. It involves referring back to an antecedent without making any explicit mention of it at all. As its name suggests, \textsc{ellipsis} involves leaving something out. But doing so, we're indicating that the elided phrase is so obvious and available that it goes without saying. And by not saying it, we paradoxically make it even more salient, as if we were tossing it up above the other semantic elements we're juggling.

Let's look at an example:

\ea \textit{Aden can \uline{speak French}, and Bella can \uline{\phantom{speak French}} too.}\\
    \begin{tikzpicture}[overlay, remember picture]
        \node[align=left,inner sep=0] (text) {\phantom{\textit{Aden can speak French, and Bella can \underline{\phantom{speak French}} too.}}};
        \draw[->,red,very thick] ([xshift=16ex,yshift=0.2ex]text.north east) to [bend left=20] ([xshift=-10.5ex,yshift=0.2ex]text.north east);
    \end{tikzpicture}
\z

Here, \textit{speak French} is missing after \textit{Bella can}, but the meaning is unmissable. Ellipsis, then, is like an invisible pronoun. It's pointing back to something, but instead of using a small word to do it, it's no word. And just like with pronouns, the antecedent needs to be available and ideally topical or nearby. Ellipsis can't happen just anywhere:

\ea[*]{\textit{This is Bella. Bella can \underline{\phantom{???}}.}} (Without prior context)
\z

Now, you might be wondering what all this has to do with our elusive independent \mention{whose}. Well, it turns out that independent \mention{whose} is doing both these tricks at once. It's acting like a pronoun, pointing back to a possessor, but it's also involved in a kind of ellipsis, leaving out the thing that's possessed. And this dual nature is part of what makes it so rare and tricky to use. But even that's not the whole story.

\bigskip

There's another kind of anaphoric relationship at play here. Let's look back at our earlier example:

\ea \textit{Aden just bought \uline{a new car}. That guy sure loves \uline{his rides}.}\\
    \begin{tikzpicture}[overlay, remember picture]
        \node[align=left,inner sep=0] (text) {\phantom{\textit{\uline{Aden} just bought a new car. \uline{That guy} sure loves cars.}}};
        \draw[->,red,very thick,dotted] ([xshift=24ex,yshift=0.5ex]text.north east) to [bend left=20] ([xshift=-4ex,yshift=0.5ex]text.north east);
    \end{tikzpicture}\label{ex:Aden2}
\z

In the previous example, Aden and that guy are the identical person. This is how anaphora usually works, with two phrases referring to the same individual or item, but here \textit{his rides} doesn't refer just to the identical car. Instead, it references a type: vehicles owned by Aden, and \textit{a new car} is just one example of this type.

Here's another example:

\ea \textit{Aden took his \uline{car} and I took mine \uline{\phantom{car}}.}\\
    \begin{tikzpicture}[overlay, remember picture]
        \node[align=left,inner sep=0] (text) {\phantom{\textit{Aden drove \uline{his car} and I drove \uline{mine}.}}};
        \draw[->,red,very thick,dotted] ([xshift=16ex,yshift=1ex]text.north east) to [bend left=20] ([xshift=-2ex,yshift=1ex]text.north east);
    \end{tikzpicture}
\z

We'll say that pronouns like \mention{my} are \textsc{dependent}, because they can't appear without a following noun, but pronouns like \mention{mine} are \textsc{independent}; they won't appear with one. (Remember that, above, we saw that there's an independent \mention{whose}. Unfortunately, while \mention{my} and \mention{mine} are easy to tell apart, \mention{whose} -- like \mention{his} -- is a bit confusing because the dependent and independent forms are spelled the same and sound the same.)

When I use the independent \mention{mine} in this example, I mean \enquote{my car}, and that meaning comes through an anaphoric relationship between the missing word and \mention{car} in the previous clause. But, again, it's not the identity-anaphora that exists between \mention{Aden} and \mention{that guy}; It's a type-anaphora. The meaning is that I took the same type of thing that Aden took -- a car -- but not the identical car that he took -- not his car.

% TODO: cut — LLM recap
We're almost at the end of this journey, but we need to pull this all together. Before we do so, let's quickly recap what we've uncovered in our linguistic investigation. We've seen that anaphora allows us to refer back to previously mentioned entities, whether through pronouns or ellipsis. We've distinguished between identity-anaphora, where we refer to the exact same entity, and type-anaphora, where we refer to the same kind of thing. We've also explored the difference between dependent pronouns like \mention{my}, which require a following noun, and independent pronouns like \mention{mine}, which stand alone. Each of these concepts plays a crucial role in understanding the behavior of our elusive independent relative \mention{whose}. Now, armed with these insights, we're ready to solve the mystery of why this particular \mention{whose} is so rare and why it appears only under very specific conditions.

\bigskip

Because \mention{mine} is a pronoun, it has an antecedent too. There's the elided \enquote{car-type} antecedent, and there's me myself. It's my car. We'll say I'm the possessor and the car is the possessed (shades of Stephen King's \textit{Christine}). So, the \textit{mine \uline{\phantom{car}}} phrase has the double anaphoric relationship shown here. The possessed car-type antecedent is related to the gap and the possessor antecedent \mention{I} is related to \mention{mine}.

\bigskip

\ea \textit{Aden took his \uline{car} and \fbox{I} took \fbox{mine} \uline{\phantom{car}}.}\\
    \begin{tikzpicture}[overlay, remember picture]
        \node[align=left,inner sep=0] (text) {\phantom{\textit{Aden drove \uline{his car} and I drove \uline{mine}.}}};
        \draw[->,red,very thick,dotted] ([xshift=19ex,yshift=0.5ex]text.north east) to [bend left=20] node[left] {possessed~~~~~~~~~~} ([xshift=-2ex,yshift=0.5ex]text.north east);
        \draw[->,red,very thick,dotted] ([xshift=14ex,yshift=5ex]text.north east) to [bend right=30] node[right] {~~~possessor} ([xshift=5.5ex,yshift=5ex]text.north east);
    \end{tikzpicture}\label{ex:Aden3}
\z

Bringing this back to independent \mention{whose}, we can now see that it also has this double anaphora. So why do we see \mention{mine} all the time, but not \mention{whose}? The thing about \mention{mine} is that, as the first person, \mention{I} is/am always topical (and so is/are \mention{you}, dear reader, as the second person). So as long as the possessed thing (or at least that type of thing) is topical, then we've met the conditions for comfortable use of \mention{mine}

But with independent relative \mention{whose}, it turns out that the possessor and possessed are rarely both topical enough at the same time. People are more likely to be topics than things, and if we're using \mention{whose}, then it's because the person is topical. But if that's the case, then the possessed usually isn't. On top of that, possessors are more likely to be people, the possessed more likely to be things, which are inherently less topical than people. Also, topicality can be spread around a little, but mostly it likes to pick out one thing. And that's fine; when the possessor is the topic, we can always use dependent \mention{whose} and avoid eliding the non-topical possessed. But it does make it a lot less surprising that independent relative \mention{whose} is so rare.

What about independent interrogative \mention{whose} you may be wondering. \textit{Whose is the brown car?} But the thing with interrogatives is that they don't \uline{refer} to antecedents -- they \uline{ask} about things. In the case of interrogative \mention{whose}, it asks about the possessor, leaving the possessed to bask in the full glow of topicality. This interrogative context presents a much less demanding topicality condition for \mention{whose} to meet.

So, let's bring this all together. The independent relative \mention{whose} that Hankamer and Postal claimed didn't exist in English does exist, but it's exceedingly rare. Its rarity is due to a burdensome confluence of linguistic conditions that need to be met:

\begin{enumerate}
    \item It needs to function as both a pronoun (referring back to a possessor) and involve ellipsis (leaving out the possessed noun).
    \item Both the possessor and the possessed need to be highly topical in the discourse.
    \item It needs to occur in a relative clause, not an interrogative context.
\end{enumerate}


When all these conditions align, we get sentences like the one from the Cambridge Grammar:

\ea{\textit{I was going to visit Lucy, }[\textit{a friend of whose had told us of the accident}]\textit{.}}
\z

Here, Lucy (the possessor) is topical because she's the person being visited, and her friend (the possessed) is topical because they're the one who told about the accident. Notice that the possessed is, unusually, a person. The oblique genitive construction \textit{a friend of whose} provides the right structural environment. And it's all happening in a relative clause.

But I've found and constructed other examples. Consider

\ea \textit{His gorilla isn't much fun at parties, but I know someone whose is.}
\z

This use of \mention{whose} has a number of factors in its favour. It's contrasting, which elevates the topicality of whatever is being contrasted. Contrasts are generally made between types -- here two members of the gorilla type -- rather than individuals, which is favourable for ellipsis. And, finally, the possessor of the gorilla isn't topical. It's just someone I know, it doesn't matter who. This allows the topical spotlight to shine on the possessed.

\section{A world tour of \mentionhead{whose}}

Our investigation of independent relative \mention{whose} in English raises an interesting question: How do other languages deal with this linguistic puzzle? Let's take a look at a few languages to see how they handle similar constructions.

We'll start with German. German has words that function similarly to our \mention{whose}: \mention{dessen} for masculine and neuter nouns, and \mention{deren} for feminine and plural. Like English \mention{whose}, these can appear without a following noun in certain contexts. Consider this example:

\ea \textit{Meins funktionierte, aber ich kenne jemanden, dessen \_ nicht funktionierte.}\\
    mine worked but I know someone whose \_ not worked\\
    `Mine was working, but I know someone whose \_ wasn't.'
\z

The underscore shows where we might expect a noun, but it's been left out. This mirrors our English construction:

\ea \textit{His gorilla isn't much fun at parties, but I know someone whose \_ is.}
\z

In both languages, we see that this construction works well in contrastive contexts, but becomes problematic in other situations. For instance, the German equivalent of our ungrammatical English sentence is equally troublesome:

% TODO: verify and fix German example — dessen + feminine Person disagrees
\ea
\gll* \textit{Die} \textit{Person}, \textit{dessen} \_ \textit{du} \textit{vergessen} \textit{hast}, \textit{ist} \textit{mein} \textit{Cousin}.\\
 {} the person whose \_ you forgotten have is my cousin\\
\glt *`The person whose \_ you forgot is my cousin.'
\z

Spanish shows a similar pattern with its word \mention{cuyo}. It's comfortable in contrastive contexts but starts to falter in constructions like this:

% TODO: verify Spanish example — cuyo without head noun
\ea
\gll* \textit{La} \textit{persona} \textit{cuyo} \_ \textit{olvidaste} \textit{es} \textit{mi} \textit{primo}.\\
   {} the person whose \_ forgot.2SG is my cousin\\
\glt *`The person whose \_ you forgot is my cousin.'
\z

% TODO: source or cut — Persian/Japanese claims need examples and citations
Interestingly, even languages that structure their sentences quite differently show analogous patterns. Persian, for example, doesn't have a direct equivalent to \mention{whose}, but it uses a possessive construction that behaves in a similar way. It's fine in contrastive contexts but becomes problematic in constructions parallel to our troublesome English examples.

Japanese takes us even further afield. It doesn't have relative pronouns at all, instead using a different strategy to build relative clauses. Yet even without a \mention{whose}-like word, Japanese still expresses these ideas in ways that echo the patterns we've seen in English and the other languages.

% TODO: cut — LLM-style numbered list
What can we learn from this linguistic world tour? A few key points stand out:

\begin{enumerate}
    \item The quirks we've observed with English \mention{whose} aren't unique to our language. Other languages show similar patterns with their possessive words in relative clauses.

    \item The preference for contrastive contexts seems to be a common thread, appearing across languages with different structures.

    \item The need for both the possessor and the possessed item to be highly salient in the conversation appears to be a shared feature.

    \item Even when languages lack a direct equivalent to \mention{whose}, they often follow similar patterns in how they express these ideas.
\end{enumerate}

These cross-linguistic patterns suggest that what we've discovered about English \mention{whose} isn't just a peculiarity of our language. It seems to tap into something more fundamental about how languages work -- how we package information, how we track what's important in a conversation, and how we balance the need for clear structure with the complexities of real-world communication.

\bigskip

It's hardly surprising that Hankamer and Postal, and many other linguists after them, thought this \mention{whose} didn't exist. Ferreting one out is like trying to spot a rare bird that only emerges under very specific conditions. If you're not looking at exactly the right time and place, you might conclude it doesn't exist at all.

But language is full of these rare constructions that push at the boundaries of what we consider grammatical.

\section{What the cluster shows}

The point isn't that friends, family, lexicographers, corpora, or language models simply outvote syntacticians. The point is that each is a different instrument. Expert judgement is one instrument, trained and sensitive, but still local. Corpus evidence is another. The \textit{OED} is another. Naive-speaker judgements and distribution-sensitive models are others. A grammaticality claim improves when those instruments are compared, not when one of them is mistaken for direct access to the language itself.

That's what the asterisk can hide. It looks like a single mark, but it may be reporting a rule, a feeling, a rarity, a processing collapse, a social boundary, or an inference from missing evidence. The missing \mention{whose} was never simply missing. It was rare, constrained, and easy to overlook. The mistake was treating that rarity as impossibility.
